\section{Module Introduction}

Artificial Intelligence has the module code COMP2261, and is led by \hyperref[modulelead-yanglong]{Dr Yang Long}.

\subsection{Assessment}


\inserttable{c|c|c|c|c|c}{Assessments for Artificial Intelligence.}
{\textbf{Title} & \textbf{Opening} & \textbf{Due} & \textbf{Feedback} & \textbf{Professor} & \textbf{Weight}}
{
    Machine Learning & 9th Oct & 15th Dec & 21st Jan & \hyperref[modulelead-yanglong]{Dr Yang Long} & $50\%$ \\
    AI Search & 15th Jan & 18th Feb & 18th Mar & \hyperref[lecturer-iainstewart]{Prof Iain Stewart} & $25\%$ \\
    Bias in AI & 5th Feb & 30th Mar & 6th May & \hyperref[lecturer-shaunaconcannon]{Dr Shauna Concannon} & $25\%$}

\subsection{Timetable}

This module is taught through two one-hour lectures and one two-hour computer class each week as seen below in \elemref{tab:Michaelmas Term Timetable for Artificial Intelligence.}{Table}.

\inserttable{c|c|c|c|c}{Michaelmas Term Timetable for Artificial Intelligence.}
{\textbf{Day} & \textbf{Time} & \textbf{Class Type} & \textbf{Location} & \textbf{Lecturer}}
{
    Monday & 14:00 - 16:00 & \colouredText{blue}{Computer Class} & TLC025 & - \\ 
    Tuesday & 16:00 - 17:00 & \colouredText{teal}{Lecture} & W103 & \hyperref[modulelead-yanglong]{Dr Yang Long} \\
    Friday & 15:00 - 16:00 & \colouredText{teal}{Lecture} & W103 & \hyperref[modulelead-yanglong]{Dr Yang Long}}

\section{Machine Learning}

This section is taught by \hyperref[modulelead-yanglong]{Dr Yang Long}. The recommended reading for this topic is: 
\begin{itemize}
    \item \hyperref[cit:Mathematics for Machine Learning]{Mathematics for Machine Learning (Deisenroth, M.P., Faisal, A.A., and Ong, C.S., 2020)}.
    \item \hyperref[cit:Mastering Machine Learning with scikit-learn]{Mastering Machine Learning with scikit-learn (Hackeling, G., 2017)}.
\end{itemize}

\subsection{Introduction to Machine Learning}

\subsubsection{Module Information}

This module is organised around a direct connection between teaching and practical work. The lectures introduce the concepts, mathematics, 
assumptions, and experimental consequences needed for that week's practical. The practical then produces the evidence used in the 
corresponding coursework report. This can be seen in \elemref{tab:The learning and coursework model.}{Table}.

\inserttable{c|c|c}{The learning and coursework model.}{\textbf{Stage} & \textbf{What you do} & \textbf{What you keep}}
{Lecture & \makecell{Understand the main question, mechanisms, \\assumptions, and key equations.} & \makecell{Notes in your own words and questions that \\still need an answer.} \\ \hline
Practical &\makecell{Implement, test, compare, visualise, and \\discuss the method.} & \makecell{A notebook that runs from top to bottom \\with outputs included.} \\ \hline
\makecell{Weekly \\Report} & \makecell{Select the strongest evidence and explain \\why the result occured.} & \makecell{A concise report that links theory, method, \\result, and interpretation.} \\ \hline
Portfolio & \makecell{Revise the nine reports and connect ideas \\across the module.} & \makecell{A coherent record of technical progress and \\developing judgement.}}

\noindent For more information about expectations for evidence, see \hyperref[cit:Machine Learning Guide]{(Long and Durham University, 2026, pp.3-4)}.

\subsubsection{What is Machine Learning?}

\definition{Machine Learning}{Arthur Samuel said in 1959 that \underline{machine learning} is the
\begin{quote}
    \say{field of study that gives computers the ability to learn without being explicitly programmed.}
\end{quote}
The field arose from the question -- can a computer learn on its own in order to perform a specific task, rather than relying on a programmer explicitly setting up 
rules for it? Similarly, Tom Mitchell defined machine learning in 1997 as
\begin{quote}
    \say{A computer program is said to learn from experience \(E\) with respect to some class of tasks \(T\) and performance measure \(P\), if its 
    performance at tasks in \(T\), as measured by \(P\), improves with experience \(E\).}
\end{quote}
From this, we can define a well-defined learning task as \(<T,P,E>\).}

\separate{7}There are three typical types of machine learning, these are:
\begin{itemize}
    \item \definition{Supervised Learning}{\underline{Supervised learning} seeks to learn the mapping (rules) between inputs and outputs. Labelled data (correct answers) is provided in past input and output pairs during the learning process to train the model on how it should behave for previously unseen data.}
    \item \definition{Unsupervised Learning}{\underline{Unsupervised learning} seeks to discover the hidden pattern from a set of inputs. Unlabelled data is provided as the past input during the learning process. The program will then attempt to group similar instances into groups.}
    \item \definition{Reinforcement Learning}{\underline{Reinforcement learning} works by providing occasional positive and negative feedback to reinforce behaviours. Good behaviours are rewarded with treats and therefore become more common whereas bad behaviours are punished and therefore become less common. This is a balance between exploration (of unchared territory) and exploitation (of current knowledge).}
\end{itemize}

\subsubsection{Differences between AI, ML, and DL}

\definition{Artificial Intelligence}{\underline{Artificial intelligence} is defined as human intelligence being exhibited by machines.}

\separate{7}\definition{Machine Learning (Simple)}{In simple terms, \underline{machine learning} can be defined as an approach to achieve artificial intelligence.}

\separate{7}\definition{Deep Learning}{\underline{Deep learning} is defined as a technique for achieving machine learning. As such, we can consider AI to be a subset of ML and ML to be a subset of DL.}

\subsection{Machine Learning Workflow and Coding Environment}

\subsubsection{Machine Learning Workflow}

\definition{Train a Model}{To \underline{train a model}, you use data to create a model and fine-tune it to make better predictions. We can then use this predictive model to make predictions on previously unseen data in order to perform specific tasks.}

\separate{7}The machine learning workflow is set out as shown below:
\begin{enumerate}
    \item \definition{Problem Framing}{\underline{Project Framing}: When setting out to create a model, we must first frame the problem. Here, we set out a clear objective that we want the model to achieve. We then define inputs and outputs and establish success metrics (e.g., prediction accuracy).}
    \item \definition{Data Preparation}{\underline{Data Preparation}: The next step is to prepare the data. To do that, we identify \underline{features} (inputs) and a \underline{label} (output). We then carry out data cleaning on the dataset and split the data into training and test sets.}
    \item \definition{Algorithm Selection}{\underline{Algorithm Selection}: Next, we select the algorithm we want to use for the model based on the dataset and the objective.}
    \item \definition{Model Training}{\underline{Model Training}: The training data is used to incrementally improve the model's performance. The progress made in this step is dependant on the quality of the work done in the previous two stages.}
    \item \definition{Model Testing}{\underline{Model Testing}: The next step is to take the model and test it against previously unseen data (the test set).}
    \item \definition{Hyperparameter Tuning}{\underline{Hyperparameter Tuning}: We can then attempt to improve the training process to further increase the model's accuracy.}
    \item \definition{Inference/Prediction}{\underline{Inference/Prediction}: Finally, the model can be deployed for the task outlined in the first step.}
\end{enumerate}

\subsubsection{Coding Environment}

When programming, we traditionally set out the features as variable \(X\) and the target label as \(Y\).

\subsection{Data Understanding}

\definition{Sample}{A \underline{sample} is the unit for which the model recieves features and, during training, a label.}

\separate{7}We then set up a feature matrix \(X\) where each row is a different sample and each column is a new feature and \(y\) is a matrix where each row is sample's label corresponding to the sample's features in \(X\).

\separate{7}\definition{Nearest-centroid classification}{\underline{Nearest-centroid classification} attempts to classify the samples by setting out one centre for each species and then classifying them based on their nearest centroid. The formula for calculating this is:
\begin{equation}
    \mu_c = \left(\frac{1}{N_c}\right) \sum_i:y^{(i)} = c x^{(i)}
\end{equation}
For each species \(c\), average the feature vectors belonging to that species. Then, subtract the centroid from the new example, square the differences, add them, then take the square root, as shown below.
\begin{equation}
    d(x,\mu_c) = ||x=\mu_x||_2
\end{equation}
To predict a sample's species, compute the distances from each species' centroid and choose the class with the smallest distance.
\begin{equation}
    \hat{y} = arg\;min_c || x - \mu_c ||_2
\end{equation}}

\subsection{Data Preparation Workflow}

\definition{Data Preparation Workflow}{To prepare data for machine learning, we use the \underline{data preparation workflow} which follows the steps as below:
\begin{enumerate}
    \item Read and inspect the data.
    \item Clean and split the data into train and test sets.
    \item Fit the baseline and model.
    \item Test and compare the model against the baseline.
    \item Scale the model and predict unseen data.
\end{enumerate}}

\separate{7}\definition{Training-set Standardisation}{We can standardise our dataset using \underline{training-set standardisation} to reduce the unbalanced impact features with large average Euclidean distances have on predictions.
To do this, we use \eqref{equation:training-set standardisation}.
\begin{equation}\label{equation:training-set standardisation}
    x' = \frac{x - \mu_{train}}{\sigma_{train}}
\end{equation}}

\subsection{Data Construction}

\subsubsection{Analysing the Raw Data}

\definition{Skewness}{Data can be considered \underline{skewed} to the right if the majority of the data lies to the left with a long tail to the right.
To measure this, we use a measure called \underline{skewness} which can be calculated with \eqref{equation:skewness}
\begin{equation}\label{equation:skewness}
    \text{Skewness} = \frac{3 \times (\text{Mean} - \text{Median})}{\text{Standard Deviation}} = \frac{3(\mu - Mdn)}{\sigma}
\end{equation}}

\separate{7}\definition{Capping}{Datasets may sometimes contain \underline{capped} features. 
This can be identified by graphing the data.
If there is a spike on the maximum-value, then it is likely that the data is capped.}

\subsubsection{Sampling}

When deciding how large of a sample to take from a population, you have to balance the size as a larger sample reduces random sampling noise, but also increases processing time.

\separate{7}\definition{Kolmogorov-Smirnov}{The \underline{Kolmogorov-Smirnov test} compares the distributions of two samples. 
It returns the maximum vertical distance between two cumulative distribution functions.
\begin{equation}\label{equation:ks test}
    D = sup_x | F_{sample}(x) - F_{population}(x)|
\end{equation}
It can then be used to calculate a p-value measuring the evidence against equal distributions.
This means that if it is above 5\%, then there is insufficient evidence of a difference in the distributions.}

\separate{7}\definition{Chi-square}{Similarly, for categorical data, you can use the \underline{Chi-square goodness of fit test} which compares the expected and observed counts of each unique value.
\begin{equation}\label{equation:chi-square}
    \chi^2 = \sum_j \frac{(O_j - E_j)^2}{E_j}
\end{equation}
where \(O_j\) is the observed count and \(E_j\) is the expected count.}

\subsubsection{Building Trustworthy Evaluation Sets}

A row should only serve one role. 
Either as a training set (to fit parameters), a validation set (to compare choices), or as a test set (to get a final accuracy estimate).

\separate{7}Random splitting preserves population proportions only in expectation. 
A finite test set can contain too many or too few examples from an important region. 
Rare strata show the largest relative variation.

\subsection{Data Transformation}

\subsubsection{Numerical Transforms}

\definition{Min-max normalisation}{In \underline{min-max normalisation}, the training minimum becomes 0 and the training maximum becomes 1.
Values between them retain their relative position in the observed training range.
\begin{equation}
    x' = \frac{x - min_{train}}{max_{train} - min_{train}}
\end{equation}}

\definition{Standardisation}{In \underline{standardisation}, values have the training mean subtracted, and then divided by the training standard deviation.
The transformed training column has mean near 0 and standard deviation near 1.
\begin{equation}
    x' = \frac{x-\mu_{train}}{\sigma_{train}}
\end{equation}}

\separate{7}\definition{Estimator Interface}{The \underline{estimator interface} has three functions: fit (to learn the statistics), transform (to apply statistics), and fit\_transform (combines on training).
It is also used in imputers, encoders, ColumnTransformer, and Pipeline.}

\separate{7}\definition{Data Leakage}{\underline{Data leakage} occurs when preprocessing learns from rows that should remain unseen.
Test-set statistics carry information about the future evaluation distribution back into training.}

\subsubsection{Skew and Bins}

To reduce the impact of skewed columns, they can be log scaled to compress multiplicative differences into smaller additive differences.
\begin{equation}
    x' = \log(1+x)
\end{equation}
The transformed distribution spreads common values while compressing the tail.

\separate{7}If you use equal-width bins, you may end up with sparse bins if data is skewed.
To fix this, we use quartile bins to keep roughly the same number of rows in each bin.

\subsubsection{Categories and Missing Values}

\definition{one-hot}{If a column is unordered, then we can apply a \underline{one-hot representation} where each unique value in the column becomes a new column with a boolean value of whether it is present.
This is done with a \underline{one-hot encoder} which learns categories during fit and then applies the same column order during transform.}

\separate{7}If you have a column with missing data, you can take one of two approaches:
\begin{itemize}
    \item Delete the incomplete rows - this does remove observations and can impact the data distribution.
    \item Impute a fitted statistic - preserves row count and creates a value for future rows (normally median).
\end{itemize}

\subsubsection{A Leakage-Free Preprocessing Pipeline}

Column-specific preprocessing uses ColumnTransformer to take each column and applies median imputation and standardisation for numeric columns, and categorical encoding for categorical ones.

\separate{7}The end-to-end data flow becomes: taking the raw DataFrame and fit preprocessing, then taking the resulting matrix and training the model.

\subsection{Machine Learning Algorithms}

\subsubsection{What a Learning Algorithm Contains}

\definition{Learning Algorithm}{A \underline{learning algorithm} consists of the five parts:
\begin{itemize}
    \item The \underline{representation} defines possible predictions. 
    \item The \underline{parameters}.
    \item The \underline{objective} scores each parameter choice.
    \item The \underline{optimiser} searches.
    \item The \underline{Hyperparameters} control the search rather than the fitted relationship.
\end{itemize}}

\noindent\definition{Representation}{A \underline{representation} states which relationships the model can express.
A straight line can change its intercept and slope.
It cannot bend, create a threshold, or model separate local patterns.}

\separate{7}\definition{Parameters/Hyperparameters}{\underline{Parameters} determine the fitted line.
The algorithm estimates them from the training data. 
Whereas, the \underline{hyperparameters} such as learning rate \(\alpha\) and iteration count control optimisation.
The algorithm does not estimate them through the same update rule.}

\separate{7}\definition{Objective Function}{The \underline{objective} converts a complete set of predictions into one number.
A lower cost means the current parameter values match the training observations more closely under the chosen definition of error.}

\separate{7}\definition{Optimiser}{An \underline{optimiser} proposes new parameter values.
Gradient descent uses local slope information from the objective.
Other solvers can reach the same least-squares solution by different routes.}

\subsubsection{A Straight-Line Model}

\definition{Simple Linear Regression}{A \underline{simple linear regression model} assigns one predicted value to each input \(x\).
The intercept sets the prediction at \(x=0\).
The slope states how much the prediction changes when \(x\) increases by one unit.
\begin{equation}
    \hat{y} = \theta_0 + \theta_1x
\end{equation}}

\separate{7}\definition{Hypothesis}{The \underline{hypothesis} is a single pair of parameters that produce predictions for every observation.
\begin{equation}
    h\theta(x_i) = \theta_0 + \theta_1x_1
\end{equation}}

\separate{7}\definition{Residual}{The \underline{residual} measures the vertical gap between prediction and observation.
\begin{equation}
    e_i = \hat{y}_i - y_i
\end{equation}}

\separate{7}\definition{Least-Squares Cost}{The \underline{least-squares cost} averages squared residuals across all \(m\) training observations.
\begin{equation}
    J(\theta_0,\theta_1) = \frac{1}{2m}\sum_{j=1}^{m} (h\theta(x_i) - y_i)^2
\end{equation}}

\separate{7}\definition{Parameter Space}{Each point in \underline{parameter space} represents one complete line. 
Moving horizontally changes \(\theta_0\) and moving vertically changes \(\theta_1\).
The height \(J(\theta_0,\theta_1)\) records how badly the line fits the data.}

\subsubsection{Gradient Descent}

\definition{Gradient Vector}{The \underline{gradient vector} is defined as follows, where both partial derivatives aare evaluated using the same current parameter values:
\begin{equation}
    \nabla J(\theta_0,\theta_1) = \left[\frac{\delta J}{\delta \theta_0}, \frac{\delta J}{\delta \theta_1}\right]^T
\end{equation}}

\separate{7}Adapting the equation, we can get the gradient with respect to the intercept.
\begin{equation}
    \frac{\delta J}{\delta \theta_0} = \frac{1}{M} \sum_i h\theta(x_i) - y_i
\end{equation}

\separate{7}Adapting the equation, we can get the gradient with respect to the slope.
\begin{equation}
    \frac{\delta J}{\delta \theta_1} = \frac{1}{M} \sum_i (h\theta(x_i) - y_i)x_i
\end{equation}

\subsubsection{The Controlled Experiment}

When viewing a learning-rate curve we can adapt the hyperparameters based on the emerging patterns.

\inserttable{c|c|c}{Reading the learning-rate curves.}{\textbf{Pattern} & \textbf{What The Curve Shows} & \textbf{Likely Response}}
{Slow Descent & Cost decreases. & Increase \(\alpha\) cautiously.\\
Smooth Convergence & Cost fails and stabilises. & Retain and verify predictions.\\
Oscillation & Cost alternates. & Reduce \(\alpha\).\\
Divergence & Cost grows rapidly. & Reduce \(\alpha\) and inspect scaling.}

\subsection{Linear Regression}

\subsubsection{Three Routes to the Same Least-Squares Line}

The iterative route uses gradient descent.
It starts from an initial guess and repeatedly reduces the cost.
Its result depends on learning rate, iteration budget, and numerical behaviour.

\separate{7}The direct route uses the closed form to calculate the least-squares parameters from summary quantities.
\textit{LinearRegression} packages a robust estimator behind a standard interface.
\begin{equation}
    \theta_0 = \frac{\sum_i (x_i - \overline{x})(y_i - \overline{y})}{\sum_i (x_i-\overline{x})^2}
\end{equation}

\subsubsection{Regression Metrics}

Mean Square Error:
\begin{equation}
    MSE = \frac{1}{m}\sum_{i} (\hat{y}_i - y_i)^2
\end{equation}

\noindent Root Mean Squared Error:
\begin{equation}
    RMSE = \sqrt{MSE} = \sqrt{\frac{1}{m}\sum_{i} (\hat{y}_i - y_i)^2}
\end{equation}

\noindent Mean Absolute Error:
\begin{equation}
    MAE = \frac{1}{m} \sum_i |\hat{y}_i - y_i|
\end{equation}

\noindent\definition{CoD}{The \underline{coefficient of determination} compares squared error with the test-set mean baseline.
\begin{equation}
    R^2 = 1 - \frac{\sum_i (y_i - \hat{y}_i)^2}{\sum_i (y_i - \overline{y})^2}
\end{equation}} 

\subsubsection{Implementation and Optimisation Variants}

\definition{Batch gradient descent}{\underline{Batch gradient descent} uses every training row for each update. 
The path is stable, but each update costs a full pass through the data.}

\separate{7}\definition{Stochastic and mini-batch}{\underline{Stochastic gradient descent} uses one row per update.
\underline{Mini-Batch gradient descent} uses a small group such as 32 rows. 
Updates become noisier but more frequent.}

\section{Artificial Intelligence Search}

This section is taught by \hyperref[lecturer-iainstewart]{Professor Iain Stewart}. There is no recommended reading for this topic however there may be an additional video to watch after lectures to enable more content to be covered. 

\separate{5}\comment{Watch the FAQs video.}

\subsection{Coding Real-World Problems for Search Tree Expansion}

\comment{Watch ahead of time - this will not be covered in lectures. Week 5 covers no content.}

\subsection{Uninformed Search: BFS, DFS, and Iterative Deepening}

\comment{Watch ahead of time - this will not be covered in lectures. Week 5 covers no content.}

\subsection{General Heuristic Search: Greedy Best-First Search and A* Search}

\subsection{A* Search and Deriving Iterative Deepening A* Search}

\subsection{Heuristic Functions and Their Development}

\subsection{Local Search: Hill-Climbing, Simulated Annealing, and Genetic Algorithms}

\subsection{Swarm Behaviour, How Ants Forage, and a Simple S-ACO algorithm}

\subsection{Applying Different Ant Systems to the TSP}

\subsection{General Versions of Particle Swarm Optimisation and Differential Evolution}

\subsection{Applying Particle Swarm Optimization to the TSP}

\section{Ethics and Bias in AI}

This section is taught by \hyperref[lecturer-shaunaconcannon]{Dr Shauna Concannon}. The recommended reading is handled on a weekly basis for this topic.

\subsection{Introduction to AI Ethics and Related Concepts} 

\subsection{Societal Impacts of AI}  

\subsection{Defining Bias and Fairness in ML}

\subsection{How to Build Better Systems}

\subsection{Value Sensitive Design}

\subsection{Evaluations and Audits} 

\subsection{Bias Mitigation Strategies}

\subsection{Transparency and Accountability for Trustworthy Systems}

\subsection{Explainable AI}

\subsection{Future Trends and Ethical Frontiers}